%=============================================================================
% MÓDULO 02 — Introdução com Fichas Críticas das Fontes Auditáveis
%=============================================================================
% Padrão: A introdução deve citar as duas fontes "ativas e auditáveis"
% fundamentais com fichas críticas completas (já apresentadas no Módulo 01).
%=============================================================================

% --- CONTEXTO DO PROBLEMA ---
O rápido crescimento de ecossistemas de IA multiagente tem levantado questões inéditas sobre autoria, integridade e responsabilidade técnica. Em sistemas como o \textbf{Core} — a orquestração central metacognitiva do OpenCode Ecosystem — as decisões não são tomadas por um único autor, mas surgem da interação de múltiplos agentes especializados, cada um com perfis de confiança, capacidades e históricos de execução (Engine Trust).

% --- A PERSPECTIVA INOVADORA ---
A perspectiva inovadora deste artigo é considerar o Core não como um sistema autônomo com "mente própria", mas como uma \textbf{prótese cognitiva humana} — uma tecnologia que estende emedia a cogniência humana, similarmente ao como um microscópio estende a visão ou uma calculadora estende o cálculo. Quando um pesquisador usa o Core para orquestrar tarefas complexas, ele está expandindo sua capacidade cognitiva através desse sistema agente.

% --- A RECONFIGURAÇÃO DO CONCEITO DE ASSINATURA ---
Esse enquadramento tem consequências diretas para o conceito de assinatura: se o Core é uma prótese, então a assinatura cognitiva não é "do sistema" isoladamente, mas da \textbf{parceria humano-prótese}. O artigo expande a definição do \textit{Prefácio R909} (SPEC-935-R909) — "assinatura cognitiva como interpretação situada de escolhas documentadas" — para o domínio da tecnologia aumentada, onde a assinatura marca como a cogniência humana se estende e se media através da arquitetura Core.

% --- FUNDAMENTAÇÃO NAS DOISFONTES AUDITÁVEIS ---
Em documentos anteriores do ecossistema, especialmente o \textit{Prefácio R909}, a assinatura cognitiva foi introduzida como "interpretação situada de escolhas documentadas, sem atribuição de personalidade ou diagnóstico universal". Este artigo situa essa definição dentro da teoria da tecnologia como prótese, onde a assinatura marca a fronteira entre a cognição nativa humana e sua extensão tecnológica.

% --- REFERÊNCIA 1 AUDITÁVEL: CLARK & CHALMERS (1998) ---
% A primeira fonte "ativa e auditável" fundamenta a metáfora da prótese cognitiva.
% Ficha crítica resumida para esta seção:
% \begin{critica}[CLARK & CHALMERS, 1998]
% \textbf{Fragmento Original:} "Quando um agente externo (por exemplo, um caderno de notas) é usado de tal forma que, se fosse removido, o agente teria dificuldade em realizar uma tarefa, então esse agente externo faz parte do sistema de pensamento do agente."
% \textbf{Relevância:} Fundamenta a ideia de que ferramentas podem ser partes integrantes do sistema de pensamento quando integradas adequadamente, legitimatizando o Core como extensão da cogniência humana.
% \end{critica}

% --- REFERÊNCIA 2 AUDITÁVEL: VERBEEK (2005) ---
% A segunda fonte "ativa e auditável" fundamenta a análise ética e da medicação tecnológica.
% Ficha crítica resumida para esta seção:
% \begin{critica}[VERBEEK, 2005]
% \textbf{Fragmento Original:} "Technologies can never be value-neutral. They are material manifestations of human intentions and, in turn, they modify the user's sense of the world and her ability to act in it."
% \textbf{Relevância:} Fornece a justificativa ética para a responsabilidade compartilhada humano-prótese, essencial para as Seções 4.1 e 4.2.
% \end{critica}

% --- JUSTIFICATIVA METODOLÓGICA ---
% O artigo adota postura \textit{fail-closed}: se um executor falha ou está ausente, informa a falha e prossegue com recursos disponíveis. Cada falha gera reflexão registrada no MetaBus e ajuste de trust. Isso estabelece um mecanismo de \textbf{accountability da prótese} em vez de culpabilização individual.

% --- QUESTÕES ÉTICAS LEVANTADAS ---
% A autoria estendida levanta questões: como citar sistemas prótese em artigos? quem responde por decisões tomadas pela prótese? como garantir que viés não seja introduzido por configurações da prótese?

% --- LIMITES METODOLÓGICOS ---
% Não invariância da prótese: a assinatura cognitiva prótese pode variar conforme atualizações de skills, mudanças no Trust Engine ou migração de agentes. Ela não é estática.

% > \textbf{Anti-overclaim:} Esta obra \textbf{não} declara que a assinatura cognitiva equivale a "personalidade do sistema" ou "consciência". Ela é interpretada como \textbf{rastreador de extensão protética}, fundamentado nas duas fontes auditáveis (Clark & Chalmers, 1998 e Verbveek, 2005). A não ser por estas duas fontes, nenhuma outra declaração sobre propriedades mentais do sistema é feita.